% ==========================================
% Copyright (c) 2026 CSharpKun & PositiveCHEEK
% This work is licensed under a Creative Commons Attribution-ShareAlike 4.0 International License.
% To view a copy of this license, visit https://creativecommons.org/licenses/by-sa/4.0
% ==========================================

\documentclass{../../libre-oge}
\usepackage{lua-ul}
\usepackage{lyluatex}
\usepackage{svg}

\author{CSharpKun \& PositiveCHEEK}
\date{2026}
\setid{00000001}
\setsubject{МУЗЫКА}
\setklas{9}

\begin{document}

\begin{center}
	\textbf{\large Инструкция по выполнению работы}
\end{center}


Экзаменационная работа состоит из двух частей, включающих в себя
19 заданий. Часть 1 содержит 14 заданий, часть 2 содержит 5 заданий с развернутым ответом.

На выполнение экзаменационной работы по ПРЕДМЕТУ отводится
3 часа (180 минут).

Ответами к заданиям части 1 является число,
последовательность цифр, слово или фраза. Ответ запишите в поле ответа в тексте работы,
а затем перенесите в бланк ответов № 1.

Решения заданий части 2 и ответы к ним запишите на бланке
ответов № 2. Задания можно выполнять в любом порядке. Текст задания
переписывать не надо, необходимо только указать его номер.

Все бланки заполняются яркими чёрными чернилами. Допускается
использование гелиевой или капиллярной ручки.

Сначала выполняйте задания части 1. Начать советуем с тех заданий,
которые вызывают у Вас меньше затруднений, затем переходите к другим
заданиям. Для экономии времени пропускайте задание, которое не удаётся
выполнить сразу, и переходите к следующему. Если у Вас останется время,
Вы сможете вернуться к пропущенным заданиям.

При выполнении части 1 все необходимые вычисления,
преобразования выполняйте в черновике. \textbf{Записи в черновике, а также
в тексте контрольных измерительных материалов не учитываются при
оценивании работы.}

Если задание содержит рисунок, то на нём непосредственно в тексте
работы можно выполнять необходимые Вам построения. Рекомендуем
внимательно читать условие и проводить проверку полученного ответа.

Баллы, полученные Вами за выполненные задания, суммируются.
Постарайтесь выполнить как можно больше заданий и набрать наибольшее
количество баллов.

После завершения работы проверьте, чтобы ответ на каждое задание
в бланках ответов № 1 и № 2 был записан под правильным номером.

\medskip

\begin{center}
	\textit{\textbf{Желаем успеха!}} 
\end{center}

\newpage

\begin{center}
\textbf{\large{Часть 1}}
\smallskip

\descbox{\textit{\textbf{%
	Ответами к заданиям 1–14 являются число, последовательность
	цифр или слово, которые следует записать в БЛАНК ОТВЕТОВ № 1 справа от
	номера соответствующего задания, начиная с первой клеточки. Если
	ответом является последовательность цифр или фраза, то запишите её \underLine[top=-0.1ex]{без
	пробелов и других дополнительных символов}. Каждый символ пишите
	в отдельной клеточке в соответствии с приведёнными в бланке
	образцами.}}
}
\end{center}

\begin{task}{1}
	Подпишите названия данных ступеней в указанном ключе.
	
	\bigskip
	
	\begin{lilypond} 
		\relative c' {
			\clef "treble" \numericTimeSignature \time 4/4 \key c \major \U c4 \U e4 \U g4
			\U e4 | % 1
			\U a4 \U g4 \U f4 \U e4 | % 2
			\U d4 \U g,4 \U a4 \U b4 | % 3
			c1 \bar "||" % 4
		}
	\end{lilypond}
	
	\answerfield
\end{task}



\begin{task}{2}
	Определите количество знаков в тональности "Ми Мажор".
	
	\answerbox
\end{task}

\begin{task}{3}
	Определите размер данной мелодии.
	
	\begin{lilypond}
		\relative c' {
			\omit Staff.TimeSignature
			\override Staff.BarLine.break-visibility = ##(#f #f #f)
			\clef "treble" 
			\time 3/4 
			\key e \minor 
			\U e8 \U g8 \U b8 \U g8 \U b8 \U g8| % 1
			\U e8 \U g8 \U b8 \U g8 \U b8 \U g8 \bar "||" % 2
		}
	\end{lilypond}
	
	\begin{numberchoices}%
		\choice 5/4%
		\choice 4/4%
		\choice 2/4%
		\choice 3/4%
	\end{numberchoices}%
	
	\answerbox
\end{task}

\begin{task}{4}
	Определите, в каких тактах содержится интервал "[Интервал]".
	
	\answerfield
\end{task}

\begin{task}{5}
	Выберите из предложенного верное название данного аккорда.
	
	\begin{lilypond}
		\relative c' 
		{
			\omit Staff.TimeSignature
			<fis a d>1 
		}
	\end{lilypond}
	
	\begin{numberchoices}%
		\choice Секстаккорд%
		\choice Квартсекстаккорд%
		\newline%
		\choice Септаккорд%
		\choice Квартсептаккорд%
	\end{numberchoices}
	
	\answerbox
\end{task}

\begin{task}{6}
	Запишите тональность данной мелодии.
	
	\answerfield
\end{task}

\begin{task}{7}
	Выпишите номера инструментов, входящих в группу меднодуховых.

	\begin{numberchoices}%
		\choice Туба%
		\choice Кларнет%
		\newline%
		\choice Тромбон%
		\choice Саксофон%
		\newline%
		\choice Гобой%
		\choice Фагот%
	\end{numberchoices}%
	
	\answerfield
\end{task}

\begin{task}{8}
	Выберите верный вариант написания Тенорового ключа.
	
	\begin{numberchoices}%
		\choice\hspace{-1.5em}\lilypond{%
			\omit Staff.TimeSignature
			\omit Staff.BarLine
			\omit Staff.ClefModifier
			\clef "treble"
			s4
		}%
		\choice\hspace{-1.5em}\lilypond{%
			\omit Staff.TimeSignature
			\omit Staff.BarLine
			\omit Staff.ClefModifier
			\clef "alto"
			s4
		}%
		\choice\hspace{-1.5em}\lilypond{%
			\omit Staff.TimeSignature
			\omit Staff.BarLine
			\omit Staff.ClefModifier
			\clef "tenor"
			s4
		}%
		\choice\hspace{-1.5em}\lilypond{%
			\omit Staff.TimeSignature
			\omit Staff.BarLine
			\omit Staff.ClefModifier
			\clef "bass"
			s4
		}%
	\end{numberchoices}
	
	\answerbox
\end{task}

\begin{task}{9}
	Определите данную длительность.
	
	\begin{lilypond}
		\omit Staff.TimeSignature
		\omit Staff.BarLine
		\relative c' {
			c256
		}
	\end{lilypond}
	
	\answerfield
\end{task}

\begin{task}{10}
	Напишите, о каком музыкальном инструменте идёт речь в тексте.
	
	\bigskip
	
	\setlength{\parindent}{25pt}
	
	Он был изобретен итальянцами в XVI в. и носил название «дульциан», что в переводе означает «нежный», «сладкозвучный». Первые версии этого инструмента были очень большими. 
	
	Его конструкция разборная и состоит из пяти съемных частей: раструба, трости, малого и большого колена и нижнего колена, которое еще называют сапогом или стволом. Корпус оснащен большим количеством отверстий, с помощью которых, попеременно открывая и закрывая их, музыкант меняет высоту звука. 
	
	Его тембр низкий и очень экспрессивный. Его партия записывается обычно в басовом либо теноровом ключах. Инструмент применяется в симфоническом, реже в духовом оркестре, а также как сольный и ансамблевый инструмент.
	
	\setlength{\parindent}{0pt}
	
	\bigskip
	
	\answerfield
\end{task}

\begin{task}{11}
	Выберите верное значение данного мелодического украшения: [Знак].
	
	\begin{lilypond}
		\relative c' {
			\omit Staff.TimeSignature
			c4\mordent
		}
	\end{lilypond}
	
	\begin{numberchoices}
		\choice Ыа
		\choice Ыа
		\choice Ыа
		\choice Ыа
	\end{numberchoices}
	
	\answerbox
\end{task}

\begin{task}{12}
	Запишите страну происхождения танца "Полька".
	
	\answerfield
\end{task}

\begin{task}{13}
	Запишите название обращенной Большой Септимы.
	
	\answerfield
\end{task}

\begin{task}{14}
	Напишите номер клавиши, обозначающей шестую ступень в тональности Фа Мажор.
	
	\includesvg[width=0.5\textwidth]{piano.svg}
	
	\answerbox
\end{task}

\begin{center}
	\warningsign
	\descbox{\textit{\textbf{%
				Не забудьте перенести все ответы в БЛАНК ОТВЕТОВ № 1
				в соответствии с инструкцией по выполнению работы.
				Проверьте, чтобы каждый ответ был записан в строке с номером
				соответствующего задания.}}
	}
\end{center}

\newpage

\begin{center}
	\textbf{\large{Часть 2}}
	\smallskip
	
	\descbox{\textit{\textbf{%
				При выполнении заданий 21–25 используйте БЛАНК ОТВЕТОВ № 2.
				Сначала укажите номер задания, а затем запишите его решение
				и ответ. Пишите чётко и разборчиво.}}
	}
\end{center}

\begin{task}{21}
	Задание с развернутым ответом
\end{task}

\begin{task}{22}
	Задание с развернутым ответом
\end{task}

\begin{task}{23}
	Задание с развернутым ответом
\end{task}

\begin{task}{24}
	Задание с развернутым ответом
\end{task}

\begin{task}{25}
	Задание с развернутым ответом
\end{task}

\begin{center}
	\warningsign
	\descbox{%
	\textit{\textbf{Проверьте, чтобы каждый ответ был записан рядом с номером
	соответствующего задания.}}
}
\end{center}

\end{document}
